% Pursuing Stacks (cotes 134-2 to 134-8): archive cote 134-3, site page 32 (PDF page 33), typescript page 32.
% typescript. Draft compared: ps2.tex:1514-1585.
% AI-transcribed by Claude Opus 5.5 (claude-opus-5-5, adaptive thinking, effort high), Message Batches API, batch ps1 (msgbatch_01Rv5hCaNTtUt4JHx1yqoHPL), 2026-11-12; image review recorded in records/CHECK_DECISIONS_134_3_4.tsv.
\scanpage{134-3/32}
\mspage{32}

generally the functor (**), admit right adjoints, namely the functor

\xout{(Cat)x\allowbreak{}x\allowbreak{}x\allowbreak{}x\allowbreak{}$\hat{A}$x\allowbreak{}x\allowbreak{}x\allowbreak{}x\allowbreak{}x\allowbreak{}x\allowbreak{}ax\allowbreak{}x\allowbreak{}x\allowbreak{}x\allowbreak{}Homx\allowbreak{}}

$j_A$: \quad (Cat) $\hand{\longrightarrow} A^{\hand{\wedge}}$ \quad , \quad where \quad $j_A(B) = (\ a \hand{\mapsto} \mathrm{Hom}(A_{/a},B))$ \quad (B $\hand{\in}$ Ob(Cat)
% NOTE: the right edge of the page cuts the end of this line; a handwritten mark (bracket?) stands there
It should be noted that the functors \quad $j_{\hand{\Delta}}$ , $j_{\hand{\square}}$ are \uline{not} the two functors \hand{$\alpha,\beta$}
% Hand-drawn arrows are marked blue; typed objects remain plain.
\diagram{\begin{tikzcd}[row sep=tiny] \& \hand{\Delta}^{\hand{\wedge}} \\ \text{(Cat)} \ar[color=blue, ur] \ar[color=blue, dr] \& \\ \& \hand{\square}^{\hand{\wedge}} \end{tikzcd}} which appear in the diagram above, \quad the latter are associated
to the familiar functors

% Hand-drawn arrows are marked blue; typed objects remain plain.
\margin{(***)} \diagram{\begin{tikzcd}[row sep=tiny] \hand{\Delta} \ar[color=blue, dr] \& \\ \& \text{(Cat)} \\ \hand{\square} \ar[color=blue, ur] \& \end{tikzcd}}

(factoring in fact through (Ord)), associating to every ordered simplex,
or to every multioriented cube, \quad the corresponding 1-skeleton with suitable
orientations on the edges, turning \inserted[typed]{the vertices of} this graph into an ordered set; \quad while
the two former are associated to the functors deduced from (**) by restric-%
ting to $A \hand{\subset} A^{\hand{\wedge}}$ , \quad namely
\[ n \hand{\mapsto} \hand{\Delta}\hand{/}\hand{\Delta}_n \]
in the case of \hand{$\Delta$}, \quad and accordingly for \hand{$\square$}. \quad The values of these functors,
contrarily to the two preceding ones, \quad are \uline{infinite} categories, \quad and they
cannot be described by \inserted{(i.e. ``are'' not)} (pre)ordered sets. \quad If however we had defined
the categories \hand{$\Delta$} , \hand{$\square$} \quad in terms of iterated boundary operations only, \quad exclu%
% NOTE: the right edge of the page cuts the end of the line after "exclu"
ding the degeneracy operations (which, \quad I feel, \quad are not really needed for
turning \hand{$\Delta^\wedge$} and \hand{$\square^\wedge$} into model\over{ato}{ize}rs), \quad we would get indeed \uline{finite} ordered
sets, \quad namely the full combinatorial simplices or cubes, \inserted[typed]{each one} embodied by the
ordered set of all it's facettes of all possible dimensions.

Contrarily to what happens with the functors \hand{$\alpha$, $\beta$}, \quad I feel that for
the two functors \hand{$\xi,\eta$} in opposite direction, \quad not only are they model pre-%
serving, \quad but the \inserted[typed]{right} adjoint functors \quad $j_{\hand{\Delta}}$ , $j_{\hand{\square}}$ must be model preserving too,
and we will have to come back upon this in a more general context.

We could amplify and unify somewhat the previous diagram of modelizers,
by introducing multicomplexes, \quad which after all can be as well ``mixed''
namely partly semisimplicial, \quad partly cubical. \quad Namely, \quad we may introduce
the would-be \inserted{``}test categories'' \hand{(30)}
\[ (\hand{\Delta})^p \hand{\times} (\hand{\Delta})^q = T_{p,q} \qquad (p,q \hand{\in} \underline{N} \ , \quad p+q \over{>}{\geq} 1) \]
giving rise to the category \quad $T^{\hand{\wedge}}_{p,q}$ \deleted{\unsure{x}} \quad of (p,q)-multicomplexes (p times
simplicial, q times cubical). \quad \xout{The} We have a natural functor (generalizing
the functors (xxx))
\[ T_{p,q} \hand{\longrightarrow} \text{(Ord)} \hand{(} \hand{\hookrightarrow} \text{(Cat)} \hand{)} \ , \]
associating to a system of p standart simplices and q standart cubes
(of variable dimensions), \quad the \uline{product} of the p+q associated ordered sets.
We get this way a functor
\[ \hand{\alpha_{p,q}} \ \deleted{\hand{\illegible{1}_{p,q}}} : \text{(Cat)} \hand{\longrightarrow} T^{\hand{\wedge}}_{p,q} \]
